% !TeX program = pdflatex
% conecthus/fundamento.tex — o LMS (Conecthus/SALCOMP) lido pelas oito leis.
% Auto-contido: compila sozinho.  pdflatex fundamento.tex
\documentclass[11pt,a4paper]{article}
\usepackage[utf8]{inputenc}\usepackage[T1]{fontenc}\usepackage[portuguese]{babel}
\usepackage{amsmath,amssymb,amsthm}
\usepackage[margin=2.3cm]{geometry}
\usepackage{booktabs}\usepackage{xcolor}
\usepackage[hidelinks]{hyperref}

\definecolor{cinza}{gray}{0.35}
\newcommand{\code}[1]{\texttt{\small #1}}
\newcommand{\medido}[1]{\par\medskip\noindent{\small\color{cinza}\textbf{Medido.} #1}\par\medskip}
\newtheorem{teorema}{Teorema}
\newtheorem{definicao}[teorema]{Definição}
\newtheorem{corolario}[teorema]{Corolário}
\newtheorem{obrigacao}[teorema]{Obrigação}
\newtheorem{observacao}[teorema]{Observação}
\newcommand{\dg}{^{\dagger}}

\title{\textbf{O Lean reverte}\\[3pt]
  {\large o sistema dirigido por LLM, lido pelas oito leis --- e onde não reverte, paga}}
\author{Aarão Melo Lopes}
\date{10 de agosto de 2026}

\begin{document}
\maketitle

\begin{abstract}
\noindent O LMS~\cite{lms} --- a plataforma Conecthus/SALCOMP que gera as nove ferramentas Lean por um
LLM com recuperação e validação em três estágios --- funciona, e o seu paper mede-o com honestidade:
nove artefactos válidos, ganhos projetados, e \emph{seis} limitações listadas sem disfarce. Este documento
não repete a medição: dá-lhe o \textbf{fundamento}. Lido pelas oito leis
(\emph{corpo estelar}, \code{sub:oitoleis}), o sistema já pratica a lei em três sítios sem a nomear ---
o validador semântico re-deriva o que o modelo afirmou e lê o \emph{resíduo}, que é o critério inteiro
da teoria («o medidor é $0$: não se compara, reverte-se»); das \emph{três} verificações que ele faz,
duas escrevem o \textbf{par soma/produto} que a cruz exige --- a soma que mede (Pareto,
$\Sigma=100\%$) e o produto que ordena (GUT $=G\times U\times T$) --- e a terceira é a legibilidade do
PDCA, o alvo numérico sem o qual a volta não tem resíduo; e a recusa deliberada do \emph{fine-tuning}
é a regra «liga-se, não se funde» aplicada antes de a conhecer. E das seis limitações que o paper
enumera, \textbf{cinco são uma}: a mesma frase da teoria --- \textbf{todo passo reverte, ou paga} ---
dita em cinco sítios onde o sistema dá um passo que não volta: a projeção sem medida, a referência que
o próprio modelo escreve, a régua que vive só no prompt, o ciclo que nunca fecha. A sexta --- a
qualidade semântica --- \emph{não} se reduz a ela: uma contramedida plausível e falsa reverte com
resíduo $0$, e o que a apanha é outro mecanismo, o resíduo na fronteira entre artefactos. É a exceção
que torna a redução falsificável. Daí a
contribuição: o \emph{roadmap} da próxima fase não precisa de ser gosto de engenharia --- \textbf{deriva-se}.
O fecho do ciclo não é RLHF nos pesos (dar estado à interface é fundir): é escrever a correção no
\emph{banco} recuperável. O medidor de um artefacto é a sua \emph{volta}: o resíduo entre o projetado e o
medido na implantação. E cada ganho grande pede o \emph{controlo} --- o mesmo objecto ao acaso, com a mesma
magnitude --- antes de se chamar resultado. Fecha-se com a fronteira \emph{decidida por mutação}, não
por gosto: o 3 do Kanban carrega estrutura; o 6 do Ishikawa e o 5 dos Porquês são livres --- mudam-se,
e nada do que desce muda. E o documento não fica na prosa: as afirmações de conta correm na bateria em
três medidores (\code{tests/lean\_cruz.c}, \code{lean\_controlo.c}, \code{lean\_convencao.c}), com as
asserções provadas por mutação.
\end{abstract}

\tableofcontents

\section{A régua herdada: as oito leis, e o critério}\label{sec:regua}

A teoria~\cite{teoria,estelar} assenta numa lei («a unidade é») e nas suas oito interpretações
operacionais --- uma por degrau, cada uma com o seu operador e o seu medidor a fechar com resíduo $0$
(\emph{corpo estelar}, \code{sub:oitoleis}):

\begin{center}\small
\begin{tabular}{@{}cll@{}}
\toprule
 & a lei & o enunciado \\
\midrule
\textbf{0} & a divisão do zero & $0=(+1)\oplus(-1)$, \ $0\dg=\infty$ \\
\textbf{1} & o dual & $1\dg=-1$ --- a involução $\nu\circ\nu=\mathrm{id}$ \\
\textbf{2} & o bidual & $K^{**}=K$ --- a volta que devolve \\
\textbf{3} & o trial & $\{-1,0,+1\}$ --- os três sinais de todo regime \\
\textbf{4} & a tetral & $T+T^{*}$ --- a cruz, órbita de quatro \\
\textbf{5} & a pental & $x^{2}=-1$ --- o bit $i$, a espiral \\
\textbf{6} & a hexal & soma $=$ produto --- a interface \\
\textbf{7} & o octonião dual & $\mathbb{H}\times\mathbb{H}^{*}$ --- ligar sem fundir \\
\bottomrule
\end{tabular}
\end{center}

\noindent E duas operações que tudo o resto instancia. A \textbf{estaca} $x\mapsto x\dg$ move, e é
involutiva. A \textbf{cruz} $x^{\times}=(x\oplus x\dg,\;x\otimes x\dg)$ não move: mede o que não se
moveu --- e tem \emph{duas} coordenadas porque uma só não chega: a \textbf{soma} é invariante ao longo
da estaca e por isso \emph{mede}; o \textbf{produto} satura no ponto fixo e por isso \emph{ordena}.

\smallskip
\noindent O critério que atravessa a teoria inteira cabe numa frase, e é a que este documento vai
aplicar ao LMS:

\begin{quote}
\textbf{O medidor é $0$.} Não se compara contra um valor esperado --- \emph{reverte-se}, e lê-se o
resíduo. E todo passo reverte, ou paga: um passo que se desfaz não apaga nada; um que só vai num
sentido apaga, e apagar é a única operação que custa~\cite{arquitetura}.
\end{quote}

\noindent Uma consequência arquitetural, herdada de~\cite{arquitetura} e usada adiante: \textbf{a
interface de transformação não tem estado}. O que transforma não guarda; o que se guarda é o
\emph{banco} --- os arquivos de quem opera ---, e os dois ligam-se por uma interface reversível
\emph{sem se fundir}. Dar estado à interface é fundir os dois, e fundir dissipa.

\section{O Lean carrega o par, e as ferramentas já o escrevem}\label{sec:lean}

O Lean define-se pela eliminação do desperdício --- e eliminar é \emph{dividir}: separa-se o que
agrega valor do que não agrega, e descarta-se metade. A teoria diz o que uma divisão custa: dividir
perde distinções, e \textbf{a dualidade é a memória da divisão} --- o que guarda a segunda metade para
a conta poder fechar~\cite{teoria}. O Lean maduro sabe-o na prática: as ferramentas que
\emph{eliminam} registam o que havia antes --- o VSM e o fluxograma escrevem os dois estados ---, e as
restantes não eliminam: ordenam, escolhem ou escrevem. As nove ferramentas do LMS, olhadas uma a uma:

\begin{center}\small
\begin{tabular}{@{}lll@{}}
\toprule
a ferramenta & o par que ela escreve & a estrutura da teoria \\
\midrule
VSM & estado \textbf{atual} $\leftrightarrow$ estado \textbf{futuro} & o dual: os dois membros, nomeados \\
Fluxograma & processo \textbf{atual} $\leftrightarrow$ \textbf{futuro}, nó a nó & o dual, outra vez --- e nomeado \\
Pareto & frequências que \textbf{somam} $100\%$ & a leitura aditiva --- \emph{mede} \\
GUT & o \textbf{produto} $G\times U\times T$ que ranqueia & a leitura multiplicativa --- \emph{ordena} \\
5 Porquês & desce da falha à \textbf{raiz} (contrai) & a erosão $\varepsilon$ --- escolhe \\
5W2H & sobe da raiz às \textbf{ações} (expande) & a dilatação $\delta$ --- escreve \\
PDCA & Plan$\to$Do$\to$Check$\to$Act$\to$Plan & o ciclo de período quatro, fechado pela leitura \\
Kanban & \emph{To Do} / \emph{In Progress} / \emph{Done} & o eixo de três estados do trânsito \\
Ishikawa & \textbf{6} categorias de causa & não escreve par: catálogo de arrumação --- e o $6$ é livre, adiante \\
\bottomrule
\end{tabular}
\end{center}

\noindent Três destes blocos merecem a frase inteira, porque não são analogia --- têm conta.

\smallskip
\noindent \textbf{Pareto e GUT escrevem o par soma/produto que a cruz exige.} O par existe porque uma
leitura só não chega: o Pareto \emph{mede} --- a soma das percentagens é invariante por construção,
$\Sigma=100\%$, e é contra essa invariância que qualquer erro aparece --- e a matriz GUT \emph{ordena}
--- o score é um produto, e um produto satura: é ranking, não medida. Usar só o Pareto é medir sem
poder priorizar; usar só o GUT é priorizar sem saber quanto pesa. O Lean tradicional usa os dois sem
dizer porque são \emph{dois}; a cruz di-lo: soma mede, produto ordena, e nenhuma sozinha é o par.
\emph{E a fronteira desta frase, dita}: a cruz completa é $(x\oplus x\dg,\;x\otimes x\dg)$ --- as duas
coordenadas carregam a involução, e nem o Pareto nem a GUT a têm. O que o Lean realiza é o \emph{par de
leituras} aditiva/multiplicativa, não o operador inteiro: estrutura, portanto --- não identidade.

\medido{nos $125$ ternos $(G,U,T)\in\{1,\dots,5\}^{3}$: subir qualquer componente sobe o produto
(o GUT ordena), e o máximo de cada classe de soma é o terno equilibrado, $\max-\min\le1$ --- o
AM--GM, inteiro; há $1206$ pares com a \emph{mesma} soma e produtos distintos (a soma sozinha não
prioriza) e $120$ com o \emph{mesmo} produto e somas distintas (o produto sozinho não pesa); e o
controlo: admitido o $0$ na escala, a monotonia quebra em $55$ casos --- a premissa $1..5$ é
medida, não adorno (\code{tests/lean\_cruz.c} \S L2--L3).}

\smallskip
\noindent \textbf{5 Porquês e 5W2H são a adjunção mórfica.} A teoria tem o par
erosão/dilatação como adjunção $\delta\dashv\varepsilon$ (\emph{Catálogo}~\cite{catalogo}):
\emph{erode-se para escolher, dilata-se para escrever}. O 5 Porquês erode --- de muitos sintomas a uma
raiz, contraindo até ao ponto que não recua mais --- e o 5W2H dilata --- da raiz escolhida ao plano de
ações, com responsável e custo. Um sem o outro é a metade: raiz sem ação não corrige, ação sem raiz
não escolhe. O LMS gera os dois, e a ordem certa (erosão primeiro) já está no seu fluxo.

\medido{a adjunção fecha: $\delta(A)\subseteq B\iff A\subseteq\varepsilon(B)$ nos $4096$ pares do
universo de $6$ bits, com a distributiva ao lado (\code{tests/toolkit.c} \S K3); e o parâmetro
dela compõe e ordena --- dilatar por $r$ e depois por $s$ é dilatar por $r{+}s$, monótono no
objecto (\code{tests/morfico\_ordem.c}).}

\smallskip
\noindent \textbf{O PDCA fecha pela leitura, e é por isso que o alvo tem de ser numérico.} O ciclo tem
quatro fases e volta à primeira --- mas o que o fecha não é dar a volta: é o \emph{Check}, a fase em que
se lê o resíduo entre o planeado e o feito. Um PDCA cujo \emph{Plan} não fixa um alvo numérico é um
ciclo cuja volta não tem resíduo legível --- roda, e não mede. A restrição que o LMS impõe no prompt
(«o PLAN tem de conter um alvo numérico»~\cite{lms}) é exactamente esta exigência: \emph{a volta
precisa de um número para o resíduo existir}.

\subsection*{A fronteira, decidida por mutação: que número cobra, e qual é livre}

Não basta dizer que o 6M «não é» a hexal nem o cinco «não é» a pental --- não afirmar é barato, e
pode-se não afirmar qualquer coisa. A fronteira decide-se com o teste da teoria, que é um só:
\textbf{muta-se o número}. Se nada do que desce muda, o número é \emph{livre} --- convenção; se algo
contável quebra, o número \emph{cobra} --- estrutura. Medido, ele diz três coisas:

\smallskip
\noindent \textbf{O 6 do Ishikawa é livre.} Re-agrupam-se as causas em $2$ a $12$ grupos --- o «seis»
vira dois, sete, doze --- e a escolha da causa raiz não se move: resíduo $0$ em todas as partições. E
o contraste mede o que protege: um seletor \emph{hierárquico} (grupo primeiro, causa depois) move com
o agrupamento em todas as mesas. Quem protege o fluxo é \textbf{escolher por causa}, e aí qualquer
$k$ serve --- que é exactamente ser convenção, agora dito com um número.

\smallskip
\noindent \textbf{O 3 do Kanban cobra.} Fundir \emph{To Do} com \emph{In Progress} colide histórias
que os três estados separam --- e colide \emph{exactamente} as que diferem só no início, contadas por
dois caminhos que concordam. O terceiro estado é o endereço do \emph{atravessar}: sem ele, o
começado-e-não-acabado não tem onde morar. E a conta prova a \emph{cardinalidade} --- o terceiro estado
carrega informação que dois não guardam; a leitura trial por cima dela (\emph{To Do} e \emph{Done}
trocados pela inversão do fluxo, \emph{In Progress} o ponto fixo) é estrutura, dita como estrutura.

\smallskip
\noindent \textbf{O 5 dos Porquês é teto, e a lei é o ponto fixo.} A raiz chega no passo $d$ da
cadeia, não no quinto; perguntar além dela é \emph{idempotente} (resíduo $0$ --- a erosão parou); e o
$5$ nem basta ($d>5$ escapa-lhe) nem é preciso ($d<5$ chega antes). A instrução certa não é «pergunte
cinco vezes»: é «pergunte \emph{até ao ponto fixo}» --- e ela não se dita no prompt, que seria a
limitação 4.5.2 outra vez: verifica-se \emph{fora} do modelo, no estágio 3 --- se o porquê $d{+}1$ não
acrescenta nada, a erosão parou.

\medido{a mutação decide os três: $40$ mesas $\times$ $20$ partições de $2$ a $12$ grupos --- a
escolha por causa move $0$ vezes, a hierárquica move em $40$ de $40$ mesas; $298$ pares de histórias
distintas --- $0$ colisões com três estados, $16$ com dois, e $16$ é também a conta fechada (diferem
só no início); $8$ cadeias de profundidade $1..8$ --- a raiz no passo $d$ em todas, idempotente além
dela, $3$ escapam ao teto de $5$ e $4$ chegam antes (\code{tests/lean\_convencao.c}).}

\section{A validação do LMS já é a volta --- pela metade}\label{sec:volta}

O terceiro estágio da validação do LMS~\cite{lms} faz uma coisa que o distingue de quase toda a
validação de saída de LLM publicada: \textbf{não compara contra um gabarito --- re-deriva}. O score GUT
é recomputado a partir dos $G$, $U$, $T$ que o próprio modelo emitiu; as percentagens do Pareto são
somadas e confrontadas com $100$; o alvo do PDCA é procurado no PLAN. É o critério da
Secção~\ref{sec:regua} em operação: reverte-se, e lê-se o resíduo.

\smallskip
\noindent E os números do paper confirmam a regra que a teoria enuncia. As restrições aritméticas
\emph{no prompt} --- a régua dita ao modelo --- deixam passar cerca de $6\%$ de artefactos com o
produto errado ou a soma fora dos $100\%$~\cite{lms}; o estágio que \emph{re-deriva} apanha todos os
\emph{desse tipo} --- os que a re-derivação torna legíveis. E a conta honesta obriga a dizer o resto:
os $6\%$ são a contagem \emph{do próprio validador}, que por construção não conta o que lhe escapa ---
a taxa de escape das classes que ele não re-deriva é desconhecida, e só uma segunda régua a mediria. A
frase da teoria para isto já estava escrita: \textbf{a régua não transporta; a volta
transporta}~\cite{teoria}. Uma régua que vive no texto do prompt não atravessa a troca de modelo, a
mudança de quantização, a borda da janela de contexto --- o próprio paper lista essa fragilidade como a
sua limitação 4.5.2. A volta atravessa, porque não depende de o modelo obedecer: depende só de a conta
fechar.

\smallskip
\noindent E há uma instância da régua que não segura \emph{à vista no próprio paper}: a restrição
ditada pede as colunas \emph{To Do} / \emph{In Progress} / \emph{Done} (Tabela~1 de~\cite{lms}), e o
quadro Kanban publicado como resultado sai com \emph{quatro} --- \emph{To Do}, \emph{In Progress},
\emph{In Validation}, \emph{Completed}, sem \emph{Done} --- e recebe validade estrutural $5/5$
(Tabela~2 de~\cite{lms}). Não é acusação: é a régua ditada a não segurar nem no caso de vitrine,
documentada pelos próprios autores.

\medido{a soma mede \emph{porque pode falhar}: em $400$ tabelas de Pareto, $215$ emissões com as
percentagens arredondadas uma a uma fecham em $100$ e $185$ falham --- e a re-derivada exata, nos
numeradores sobre o denominador comum, fecha em todas, sem uma divisão; o denominador
desatualizado deriva \emph{exactamente} $100\,(F'{-}F)$ --- a lei, não uma tolerância. E a volta
contra a régua: $15$ erros injetados nos $125$ ternos GUT --- a re-derivação apanha os $15$, a
régua de alcance $1..125$ (o que uma restrição de esquema vê) apanha $0$
(\code{tests/lean\_cruz.c} \S L1, \S L4).}

\begin{observacao}[a validação interna é verdadeira, e parcial]\label{obs:parcial}
O validador do LMS lê o resíduo do artefacto \emph{contra si próprio}: a aritmética do modelo contra
os operandos do modelo. Esse resíduo é metade de um par. A outra metade é o resíduo do artefacto
\emph{contra o mundo}: o lead time projetado contra o medido na estação, a redução prometida contra a
obtida. O sistema atual lê a primeira e não tem como ler a segunda --- e um resultado que cai só de um
lado de um par dual está pela metade.
\end{observacao}

\noindent A metade que falta tem nome na teoria: \textbf{a referência escrita à mão}. Quando os números
do VSM são inferidos pelo modelo a partir da descrição textual (limitação 4.5.3 do paper), o modelo
está a escrever a referência contra a qual será julgado --- e uma referência que o próprio réu redige
valida-se sempre. O teste barato é a \emph{mutação}: mude-se o dado de entrada; se a referência não
muda sozinha, é cópia, não medida. O LMS tem hoje uma única leitura externa --- os benchmarks
documentados via RAG, contra os quais $68\%$ das estimativas caem a $\pm15\%$~\cite{lms} --- e essa
é, precisamente, a leitura a alargar.

\section{As seis limitações: cinco são uma, e a sexta muda de régua}\label{sec:seis}

O paper do LMS lista seis limitações (4.5.1--4.5.6), e a lista é honesta. Lida pelas leis, ela
\emph{colapsa} --- mas não inteira, e a exceção importa. \textbf{Cinco} das seis são a mesma frase ---
\emph{um passo que só vai num sentido apaga} --- dita em cinco sítios do sistema. A sexta, a 4.5.5, não
é: uma contramedida plausível e logicamente errada é internamente coerente --- reverte-se, e o resíduo
é $0$ \emph{sendo ela falsa}. A reversibilidade não alcança a qualidade semântica; o que a alcança é o
resíduo \emph{entre} artefactos, na fronteira (Obrigação~\ref{ob:sete}). \emph{E é esta exceção que
torna a redução falsificável}: se toda limitação coubesse na frase, a frase não diria nada.

\begin{center}\small
\begin{tabular}{@{}p{4.2cm}p{4.9cm}p{5.2cm}@{}}
\toprule
a limitação do paper & o passo que não volta & a volta que falta \\
\midrule
4.5.1 cenário único & a projeção sai sem o seu par & o par otimista $\leftrightarrow$ pessimista, nomeado --- e o \emph{controlo} (\S\ref{sec:deriva}) \\
4.5.2 lógica só no prompt & a régua não transporta & re-derivar fora do modelo --- o estágio 3, alargado \\
4.5.3 números inferidos, não medidos & a referência escrita pelo réu & a mutação do dado; a medida real na estação \\
4.5.4 sem ciclo de retorno & gera-se, e nada regressa & o resíduo da implantação, escrito no banco \\
4.5.5 validação semântica limitada & \textbf{não reduz}: o falso coerente reverte com resíduo $0$ & a fronteira entre artefactos (Obr.~\ref{ob:sete}), não a volta \\
4.5.6 aritmética do LLM ($\sim6\%$) & o enunciado sem a conta & já resolvida onde a volta existe --- a regra geral \\
\bottomrule
\end{tabular}
\end{center}

\noindent A linha 4.5.6 é a prova interna: \emph{no único sítio onde o LMS já reverte, a limitação já
não é limitação} --- o erro aritmético ocorre, é apanhado, e o sistema segue. As quatro da mesma
família (4.5.1--4.5.4) são o mesmo movimento por fazer; a 4.5.5 espera pela fronteira, que é doutra
régua. Não são seis defeitos independentes a atacar um a um: são \emph{uma} propriedade --- a
reversibilidade --- a estender sítio a sítio, mais a exceção que a delimita.

\smallskip
\noindent E o enquadramento tem borda, dita: a latência ($3{,}2$--$11{,}7$~s por ferramenta, até
$38$~s sob carga~\cite{lms}) e o custo de inferência não são passos que não voltam --- são preço de
operação, e nenhuma volta os paga. O que não cabe na frase fica fora dela.

\section{A derivação: o que a lei obriga na próxima fase}\label{sec:deriva}

O paper propõe um \emph{roadmap} (RLHF, multi-agente, integração ERP/MES, motor de cenários). A teoria
permite algo mais forte do que propor: \textbf{derivar} --- dizer qual forma cada extensão \emph{tem}
de ter, e qual forma a lei proíbe.

\begin{obrigacao}[o fecho do ciclo escreve no banco, não nos pesos]\label{ob:banco}
A interface de transformação não tem estado (\S\ref{sec:regua}); dar-lho é fundir o que devia
ligar-se. O RLHF contínuo sobre correções de especialistas é exactamente isso: estado empurrado para
dentro da interface --- com o custo conhecido de esquecer o que sabia e de não se poder auditar o que
aprendeu. A forma que a lei permite já está no sistema: \textbf{o corpus recuperável é o banco}. A
correção do especialista, o desfecho da implantação, a avaliação do artefacto --- escrevem-se como
documentos indexados, e a recuperação trá-los ao próximo prompt. O modelo continua sem estado; o banco
lembra; a estrela liga os dois sem os fundir. E note-se: o LMS \emph{já} fez esta escolha uma vez, ao
recusar o fine-tuning a favor do RAG~\cite{lms} --- a próxima fase só precisa de a repetir, agora
sabendo porquê.
\end{obrigacao}

\begin{obrigacao}[o medidor de um artefacto é a sua volta]\label{ob:residuo}
Os $65\%$ de redução de lead time e os $200\%$ de produtividade do caso da estação 04 são
\emph{projeções} --- o paper di-lo. A lei diz o que os promove a resultado: a \textbf{volta}. Cada
artefacto implantado ganha uma coluna: o valor projetado, o valor medido, o resíduo. A integração
ERP/MES do roadmap não é «mais dados para o prompt»: é \emph{o outro membro do par} da
Observação~\ref{obs:parcial} --- a leitura externa sem a qual a validação fica pela metade. E o resíduo
lido fecha as duas pontas de uma vez: valida o artefacto, e vira documento do banco
(Obrigação~\ref{ob:banco}) --- o ciclo de retorno não é um subsistema novo, é a volta a ser escrita
onde a leitura já mora.
\end{obrigacao}

\begin{obrigacao}[todo ganho grande pede o controlo]\label{ob:controlo}
Quando um número melhora muito, gera-se \emph{o mesmo objecto ao acaso, com a mesma magnitude}, e
mede-se: se o acaso empata, a métrica é degenerada --- o ganho veio do grau de liberdade, não da
escolha~\cite{teoria}. Para o LMS: ao lado do plano de contramedidas gerado, um plano \emph{aleatório}
de mesmo custo total e mesma cardinalidade, avaliado pela mesma régua. O critério é o do medidor, sem
limiar: se o acaso \emph{empata} --- empate exacto, sorteio após sorteio ---, a projeção não distingue
a escolha do modelo de escolha nenhuma, e sabê-lo \emph{antes} do investimento é o próprio serviço que
se pede a uma priorização. O motor de cenários do roadmap
(otimista/realista/pessimista) ganha assim o quarto cenário, que é o que falta aos três: o controlo.
\end{obrigacao}

\medido{o mecanismo do controlo, em miniatura: $200$ mesas de $12$ estações, $40$ planos ao acaso
por mesa --- $8000$ controlos. A régua degenerada é um KPI real e cego à escolha --- o \emph{número
de contramedidas do plano} --- e empata \emph{por construção}: essa metade não prova nada, e
di-lo-se. O que se mede é o resto: a métrica informativa (o tempo poupado, somado) \emph{discorda}
do acaso em todas as mesas, o plano do modelo não perde para sorteio nenhum, e o \emph{empate
universal} denuncia a régua constante sem ninguém ler o código dela --- é isso que o controlo dá de
graça (\code{tests/lean\_controlo.c} \S N1).}

\begin{obrigacao}[o multi-agente liga sem fundir]\label{ob:sete}
A orquestração multi-agente --- e o alvo é dos próprios autores: o roadmap deles já pede que as
conclusões do Pareto se reflitam na ordem do 5W2H~\cite{lms} --- tem duas formas possíveis, e é aqui
que a lei decide o que o roadmap deixa em aberto. A forma \emph{proibida} é a fusão: um contexto
único onde todos os artefactos se geram juntos --- aí a inconsistência entre ferramentas deixa de ser
\emph{legível}, porque não há fronteira onde a medir. A forma \emph{obrigada} é a da Lei~7 ---
\textbf{dois corpos e a interface entre eles}, $\mathbb{H}\times\mathbb{H}^{*}$, ligados sem fundir:
cada agente gera o seu artefacto, e a consistência mede-se \emph{na interface}, como resíduo contável
--- as inversões entre o ranking do GUT e a ordem do 5W2H; a raiz do 5 Porquês contida ou não num ramo
do Ishikawa; o alvo do PDCA igual ou não ao futuro do VSM. O que a lei acrescenta ao roadmap não é o
alvo --- é a forma e a métrica: \textbf{dois caminhos que têm de concordar apanham o que nenhuma
asserção interna apanha}~\cite{teoria}.
\end{obrigacao}

\medido{o mecanismo da fronteira, em miniatura: $300$ pares GUT/5W2H, com $95$ trocas injetadas na
ordem das ações (só entre scores \emph{distintos} --- trocar empatados é mutação equivalente). O
recall de $95$ de $95$ é \emph{por construção} --- a troca adjacente de valores distintos cria a
inversão que o resíduo conta --- e essa metade não se reivindica. O que se mede é o resto: a
validação \emph{interna} de cada artefacto apanha $0$ (e não por estar morta --- mutada a troca em
duplicação, acusa), e o resíduo na fronteira dá $0$ falsos nos pares consistentes. O que nenhum
artefacto vê sozinho, a fronteira conta (\code{tests/lean\_controlo.c} \S N2).}

\section{Limites: o que está afirmado, e com que régua}\label{sec:limites}

\begin{itemize}\itemsep3pt
\item \textbf{Por identidade de conta}: o validador de estágio 3 do LMS é o critério do resíduo em
  operação --- re-deriva e lê, não compara (\S\ref{sec:volta}); e cinco das seis limitações do paper
  instanciam a mesma propriedade em falta, a reversibilidade, com a 4.5.5 como a exceção que delimita
  a redução (\S\ref{sec:seis}).
\item \textbf{Por estrutura}: Pareto e GUT escrevem o par de leituras aditiva/multiplicativa que a
  cruz exige --- sem a involução, e por isso estrutura e não identidade; 5~Porquês/5W2H realizam o par
  erosão/dilatação; o VSM e o fluxograma escrevem o dual com os dois membros nomeados; o PDCA fecha
  pela leitura e por isso exige o alvo numérico. Correspondências de forma, com a conta ao lado onde
  ela existe.
\item \textbf{Por mecanismo medido}: o empate do acaso denuncia a régua cega à escolha, e o resíduo na
  fronteira conta o que a validação interna não vê --- em miniatura, com as metades triviais marcadas
  como construção e as asserções provadas por mutação (\code{tests/lean\_controlo.c}).
\item \textbf{Por mutação}: o 6 do Ishikawa e o 5 dos Porquês são \emph{livres} --- mudam-se e nada do
  que desce muda --- e o 3 do Kanban \emph{cobra}: fundir dois estados colide histórias contáveis. A
  pergunta «é a hexal? é a pental?» deixou de ser opinião: mediu-se qual número carrega estrutura e
  qual não (\code{tests/lean\_convencao.c}).
\item E a medição que falta está \emph{nomeada}, não escondida num «não afirmo»: o resíduo projetado
  contra medido, na estação --- é ela que decide o efeito das obrigações da \S\ref{sec:deriva} nos
  números do LMS, e é a régua que este documento deixa apontada.
\end{itemize}

\smallskip
\noindent \textbf{O que fica}: o LMS não precisa de importar uma teoria --- precisa de nomear a que já
pratica. Onde ele reverte, funciona e o próprio paper o mede; onde dá passos sem volta, o próprio
paper o lista como limitação. A próxima fase inteira cabe numa frase, e ela já estava
escrita~\cite{arquitetura}: \emph{todo passo reverte, ou paga.}

\begin{thebibliography}{9}
\bibitem{lms} D.~N.~Lopes, L.~F.~O.~Alves, R.~H.~L.~Prado, M.~R.~da~Cunha, R.~M.~Teixeira~Junior,
F.~C.~Novélo. \emph{LLM-Driven Lean Manufacturing System: Integrating AI, RAG, and Embeddings for
Industry 4.0 Decision Support}. Journal of Computer Science and Technology Studies 8(8): 91--102, 2026.
\code{doi:10.32996/jcsts.2026.8.8.8}.
\bibitem{teoria} A.~M.~Lopes. \emph{A teoria} (\code{teoria.tex}) --- a Lei e as suas interpretações;
a estaca, a cruz, e o critério do resíduo.
\bibitem{estelar} A.~M.~Lopes. \emph{O corpo estelar completo} (\code{papers/corpo-estelar.tex}) ---
as oito leis (\code{sub:oitoleis}) e o teorema dos tecidos.
\bibitem{catalogo} A.~M.~Lopes. \emph{O catálogo} (\code{catalogo.tex}) --- as seis leis
(\code{obs:seis-leis}); o mórfico como adjunção $\delta\dashv\varepsilon$.
\bibitem{arquitetura} A.~M.~Lopes. \emph{A arquitetura} (\code{papers/arquitetura.tex}) --- a
exigência «todo passo reverte, ou paga»; a interface sem estado e o banco.
\end{thebibliography}

\vfill
{\footnotesize\color{cinza}\noindent
Teoria, catálogo, medidores e o resto do sistema em
\href{https://goldenkingdom.patriatechnology.com}{goldenkingdom.patriatechnology.com}.
CC BY 4.0 (textos) $\cdot$ MIT (código).\par}

\end{document}
